\documentclass{ximera}
\title{Haning Chains and the Number $e$}

\newcommand{\pskip}{\vskip 0.1 in}

\begin{document}
\begin{abstract}
Catenaries the number $e$.
\end{abstract}
\maketitle


The number $e$ comes up in unexpected places. Here's one.

\section{The Hyperbolic Cosine and Sine Functions}

The function \emph{hyperbolic cosine} is defined as
\[
    \cosh(x) = 0.5(e^x + e^{-x}  ).
\]

Its graph is shown below.

\begin{onlineOnly}
    \begin{center}
\desmos{nubdwzgsa9}{450}{600}  
\end{center}
\end{onlineOnly}

\href{https://www.desmos.com/calculator/nubdwzgsa9}{152: Hyp Cosine}

\begin{enumerate}
\item Drag the slider $u$ in Line 2 of the worksheet above. Explain how to get the graph of the function $y=\cosh x$ from the graphs of the functions $y=e^x$ and $y=e^{-x}$.

\begin{freeResponse}
\end{freeResponse}

\item What is the range of the function $f(x)=\cosh x$?

\item Is the function $f(x)=\cosh x$ one-to-one? If not, what would be one way to restrict its domain to get a new function that is one-to-one and has the same range as $f$?

\item Activate the folder \emph{Hyperbolic Sine Function} in Line 15 above to see the graph of the hyperbolic sine function
\[
   \sinh(x) = 0.5(e^x - e^{-x}  ).
\]
\item Drag the slider $u$ in Line 2 of the worksheet above. Explain how to get the graph of the function $y=\cosh x$ from the graphs of the functions $y=e^x$ and $y=e^{-x}$.
\begin{freeResponse}
\end{freeResponse}

\item What is the range of the function $g(x)=\sinh x$?

\item Is the function $g(x)=\sinh x$ one-to-one?

\item Prove that 
\[
      \cosh^2x - \sinh^2x = 1.
\]

\emph{Hint: For a quick way, first factor the above expression.}


\end{enumerate}


\section{Hanging Chains}

Hold the ends of a chain and let the chain sag under its own weight. The resulting curve looks like a parabola, and so it was thought. But it turns out that this is not correct. 


\begin{onlineOnly}
    \begin{center}
\desmos{tj3dz2cnf0}{450}{600}  
\end{center}
\end{onlineOnly}

\href{https://www.desmos.com/calculator/tj3dz2cnf0}{152: Hanging Chain 1}


\begin{enumerate} 

\item To see that the dashed (orange) chain in the worksheet above is \emph{not a parabola} drag the slider $c$ in Line 2.

\item It turns out that a chain with uniform density hanging in a uniform graviatational field assumes the shape of the hyperbolic cosine function $y=\cosh x$ (we'll see why later in the course). All that's needed to describe a particular chain is to scale the size of the graph by some factor. Test this by dragging the slider $a$ in Line 6 to make the graph of the function
\[
     f(x) = -a + a\cosh(x/a)
\]
match the highlighted (orange) chain below.
\end{enumerate}


\section{Some Questions}

\begin{question} \label{QLfeenvd}
\begin{enumerate}
\item Find an expression for the inverse of the function
\[
     f(x) = \cosh x = \frac{1}{2}\left( e^x + e^{-x} \right) \, , \, x\geq 0.
\]
Start by making the substutition $u=e^x$. Solve the resulting quadratic equation by completing the square.

\item Find an expression for the inverse of the function
\[
    g(x) = \sinh x = \frac{1}{2}\left( e^x + e^{-x} \right) \, , \, x\in \mathbb{R}.
\]
Start by making the substutition $u=e^x$. Solve the resulting quadratic equation by completing the square.

\item Simplify the composition
\[
       f(g^{-1}(x))
\]
as much as possible.
\end{enumerate}
\end{question}

\begin{question} \label{Qmdfmzmcp023}
The function
\[
    h = f(x) = h_0 + a\cosh(x/a) \, , \, x\geq 0,
\]
expresses the height (in feet) of a point on hanging chain in terms of its distance (in feet) from the chain's axis of symmetry.

The function
\[
      s = g(x) = a\sinh(x/a) \, , \, x\geq 0,
\]
gives the length of the chain (in feet) measured from its low point to a point on the chain $x$ feet from the chain's axis of symmetry.

\begin{enumerate}

\item What are the units of $h_0$? How do you know?

\item What are the units of $a$? How do you know?

\item Find a function
\[
     h = w(s) \, , \, s\geq 0,
\]
that expresses the height (measured in feet) of a point $P$ on the chain in terms of the length of the chain from its low point to $P$.

\end{enumerate}

\begin{onlineOnly}
    \begin{center}
\desmos{5oppqdy2l1}{450}{600}  
\end{center}
\end{onlineOnly}

\href{https://www.desmos.com/calculator/5oppqdy2l1}{141: Catenary}


\end{question}


\end{document}